\documentclass[10pt,a4paper]{amsart}
\usepackage[margin=19mm]{geometry}
\usepackage{amsmath,amssymb,mathtools}
\usepackage[scr=rsfs,cal=euler]{mathalfa}
\usepackage{xfrac}
\usepackage[unicode,hidelinks,bookmarks=false]{hyperref}
\newcommand{\ie}{i.e.\ }
\newcommand{\cf}{cf.\ }
\newcommand{\resp}{respectively\ }
\newcommand{\Subsection}{\subsection}
\providecommand{\nobreakdash}{\nobreak\hbox{-}}
\setlength{\emergencystretch}{2em}
\multlinegap0pt
\usepackage{amssymb}
\newtheorem{prop}{Proposition}
\renewcommand{\P}{\mathcal{P}}
\title{Quasicontinuity of $N^{1,\infty}$ functions and the Vitali-Carath\'eodory property on general metric spaces}
\author{Anders Bj\"orn, Jana Bj\"orn}
\date{Source: arXiv:2605.22674v1 (CC BY 4.0)}
\begin{document}
\maketitle

\noindent\emph{Source and licence.} Quasicontinuity of $N^{1,\infty}$ functions and the Vitali-Carath\'eodory property on general metric spaces. Version arXiv:2605.22674v1 (CC BY 4.0), distributed by arXiv under the Creative Commons Attribution 4.0 International licence, http://creativecommons.org/licenses/by/4.0/, as recorded in the arXiv licence metadata for this exact version and shown on the abstract page https://arxiv.org/abs/2605.22674v1. Full licence notice: \texttt{licenses/arxiv-2605.22674-CC-BY-4.0.txt}.\par\smallskip

\noindent\textit{Editorial scope.} Counted unit: the article's prop prop-VC (source0.tex lines 545--550). The article's complete original proof of it is delivered with it (source0.tex lines 731--745, one \texttt{proof} environment of 15 lines, which the article places away from the statement). No mathematics is written, repaired or re-derived: the statement and the proof are the article's own lines, in the article's own notation, and no retained line is substituted or re-flowed.  No further source-local context is needed: every cross-reference of every retained line resolves inside the two delivered spans.  Nothing was omitted from any retained span, so the package carries no omission record.  No retained line contains a citation, so the wrapper carries no bibliography and no bibliography entry.  No retained line contains a figure environment, an image-inclusion command or drawing code, so no image is delivered and the package declares no asset dependency.  Editorial formatting only, in addition: an AMS \texttt{amsart} preamble that restates the article's own declarations and macros used by the retained lines, namely the environments \texttt{prop} on separate counters, together with the standard packages those lines need.  The retained environments print in this document as Proposition 1 in order of appearance, whereas the article prints the same items with the numbering of its own full document.  The complete native source of the article as distributed in the arXiv e-print for this exact version is bundled unchanged as \texttt{sources/201150/source0.tex}.
\begin{prop}  \label{prop-VC}
The Vitali--Carath\'eodory property for $L^\infty(\P)$ holds if and only if
$\mu(\{x\})>0$ for every $x\in\P$.

In particular, $\P$ has to be finite or countable.
\end{prop}

\begin{proof}[Proof of Proposition~\ref{prop-VC}]
If $\mu(\{x\})=0$, then $\|\chi_{\{x\}}\|_{L^\infty(\P)} =0$.
Let $v\ge \chi_{\{x\}}$ be a lower semicontinuous function.
Then $v>\tfrac12$  in an open set $G\ni x$ (which has positive measure) and hence 
$\|v\|_{L^\infty(\P)}> \tfrac12$.
This shows that the Vitali--Carath\'eodory property fails.

Conversely, assume that $\mu(\{x\})>0$ for every $x\in\P$ and $u\in L^\infty(\P)$.
Then, for every $u\in  L^\infty(\P)$, the constant function
\[
v \equiv \|u\|_{L^\infty(\P)} = \sup_{x\in \P} |u(x)|
\]
is a lower semicontinuous 
majorant of $|u|$ with $\|v\|_{L^\infty(\P)} = \|u\|_{L^\infty(\P)}$.
\end{proof}

\end{document}
